\documentclass{article}
\usepackage[T1]{fontenc}
\usepackage{lmodern}
\usepackage{graphicx}
\usepackage{amsmath,amssymb}
\usepackage[a4paper,margin=2cm]{geometry}
\usepackage[absolute,overlay]{textpos}
\setlength{\TPHorizModule}{1mm}
\setlength{\TPVertModule}{1mm}
\usepackage[indonesian]{babel}
\usepackage{hyperref}
\setlength{\emergencystretch}{2em}
% unit: Halaman 9

\begin{document}

% === Halaman 9 (python/native) ===

• Dengan panjang kunci 128-bit, maka terdapat  sebanyak


2128 = 3,4  1038  kemungkinan kunci. 

• Jika komputer tercepat dapat mencoba 1 juta kunci setiap detik, maka 

akan dibutuhkan waktu 5,4  1024 tahun untuk mencoba seluruh kunci. 

• Jika komputer tercepat yang dapat mencoba 1 juta kunci setiap 

milidetik, maka dibutuhkan waktu 5,4  1018 tahun untuk mencoba 

seluruh kunci. 

• Artinya, AES diharapkan  tahan terhadap serangan exhaustive key 

search attack (brute force attack)

\end{document}
